%% SCRAP: experiments/02-experiments/heartbeat-doe/phase2-corrected
%% SOURCE: docs/working/experiments/02-experiments/heartbeat-doe/phase2-corrected.md
%% STATUS: WORKING
%% FITS: experiments/ch-heartbeat
%% EDITORIAL: lifted — prose rewritten to press voice

\section{Heartbeat DoE Phase~2 — Corrected and Finalised}
\label{sec:heartbeat-doe-phase2-corrected}

\subsection{Executive Summary}

The initial Phase~2 design measured jitter in a fixed heartbeat interval.
The corrected design measures \emph{adaptive heartbeat rate modulation}: how
well each configuration's heartbeat rate responds to workload changes. This
section is the executive-level statement of the corrected design.

\subsection{The Physics Story}

The heartbeat is not a fixed timer; it is a feedback mechanism. The inference
engine adjusts \texttt{tick\_ns} in response to workload intensity, slowing
the heartbeat under high load to allocate more CPU time for physics-parameter
tuning:

\[
  \text{high load} \;\Rightarrow\; \uparrow\texttt{tick\_ns}
  \;\Rightarrow\; \text{longer physics tuning window}
  \;\Rightarrow\; \text{convergence to optimal rate}
\]

The primary discriminating metric for configuration selection is therefore the
Pearson correlation between per-tick workload intensity $W_t$ and per-tick
heartrate $H_t = 10^9\,/\,\texttt{tick\_ns}_t$:

\begin{equation}
  \rho = \frac{\operatorname{Cov}(W, H)}
              {\sqrt{\operatorname{Var}(W)\,\operatorname{Var}(H)}}
\end{equation}

Target: $\rho \geq 0.75$ for a configuration to qualify as a golden-configuration
candidate.

\subsection{How the Initial Design Was Wrong}

\begin{center}
\begin{tabular}{lp{5.5cm}p{5.5cm}}
\toprule
Aspect & Initial approach & Corrected approach \\
\midrule
Primary metric & Jitter in assumed-fixed 1~ms interval & Load-to-heartrate correlation \\
Thread enabled & No (\texttt{HEARTBEAT\_THREAD\_ENABLED=0}) & Yes (adaptive behaviour active) \\
Measures & Steady-state smoothness & Adaptive-response quality \\
Golden criterion & Lowest CV & Highest $\rho$ + fastest convergence \\
\bottomrule
\end{tabular}
\end{center}

\subsection{Expected Discrimination}

The corrected design is expected to reveal clear ranking separation on
$\rho$. A configuration with $\rho < 0.2$ does not adapt its heartrate to
load; regardless of static performance, it is not the golden configuration.

The expected ordering for the five elite configurations (subject to
experimental validation):

\begin{enumerate}
  \item \texttt{1\_0\_1\_1\_1\_0} — balanced loops, expected strongest coupling
  \item \texttt{1\_0\_1\_1\_1\_1} — adds decay inference, expected similar coupling
  \item \texttt{1\_1\_0\_1\_1\_1} — skips linear decay, alternative adaptation path
  \item \texttt{1\_0\_1\_0\_1\_0} — lean configuration, expected slower adaptation
  \item \texttt{0\_1\_1\_0\_1\_1} — no heat tracking, expected weakest coupling
\end{enumerate}

\subsection{Implementation Checklist}

The following code changes are required before the experiment can be executed:

\begin{itemize}
  \item Add \texttt{HeartbeatRateSample} struct to \texttt{include/vm.h}.
  \item Add \texttt{heartbeat\_rate\_samples} buffer and sample count to the
        \texttt{VM} struct.
  \item Track \texttt{workload\_ops\_current\_tick} in the execution loop;
        reset it each heartbeat cycle.
  \item Call \texttt{capture\_heartbeat\_rate\_sample()} in
        \texttt{heartbeat\_thread\_main()} before each sleep.
  \item Extend \texttt{DoeMetrics} with heartrate fields.
  \item Implement \texttt{analyze\_heartbeat\_rate\_trajectory()} in
        \texttt{src/doe\_metrics.c}.
  \item Update CSV header and row writers.
\end{itemize}

\subsection{Analysis Focus}

The R analysis script for the corrected design ranks configurations
primarily by \texttt{load\_to\_heartrate\_correlation}, secondarily by
\texttt{heartrate\_responsiveness\_score}:

\begin{lstlisting}[language=R]
ranking <- doe_data %>%
  group_by(configuration) %>%
  summarize(
    mean_correlation    = mean(load_to_heartrate_correlation),
    mean_responsiveness = mean(heartrate_responsiveness_score)
  ) %>%
  mutate(
    golden_score = (mean_correlation * 0.6) + (mean_responsiveness * 0.4)
  ) %>%
  arrange(-golden_score)
\end{lstlisting}

\subsection{Success Criteria}

\begin{itemize}
  \item Heartbeat rate samples captured per run (sample count $> 0$).
  \item CSV contains \texttt{load\_to\_heartrate\_correlation} column with
        non-zero values.
  \item Correlation metrics span a range of at least 0.4 across the five
        configurations (sufficient discrimination).
  \item Golden configuration identified with $p < 0.05$ versus second-ranked.
  \item Golden configuration achieves $\rho \geq 0.75$.
\end{itemize}

%% TODO(bob): update status to HISTORICAL once experiment is executed and results archived

